\documentclass[12pt, draftcls, onecolumn]{IEEEtran}
\usepackage{amsmath}
\usepackage{graphicx}
\usepackage{tabularx} % Put this in your preamble
\usepackage{array}    % For column formatting
\usepackage[utf8]{inputenc}
\usepackage{chemarr}
\usepackage{float}
\usepackage{cleveref}
\usepackage{amssymb}
\usepackage{booktabs}

\title{Virtual Clinical Trial Design for Infective Endocarditis: Evaluating Antibiotic Therapy}
\author{Jacob W. Kammer, Klas Udekwu, and Esteban A. Hernandez-Vargas}
\date{\today} % Use \date{} to leave it blank or specify a custom date

\begin{document}


\maketitle
 


\section{Introduction}
Infective endocarditis is a disease of the endocardial surface of the heart that involves the cardiac valves and the mural endothelium \cite{Li2024}. It is an indolent infection with Staphylococcus aureus (S. aureus) and Viridans streptococci as the most common pathogens \cite{HammondHaley2023}. It is an aggressive infection with an in hospital mortality rate of 15\% to 20\% and a one-year mortality rate of 40\% \cite{Taduru2023}. Complications include congestive heart failure, thromboembolic events, perivalvular and annular abscesses \cite{murdoch2009infective}, sepsis and septic shock. Methicillin-resistant Staphylococcus aureus (MRSA) is the most common cause of health care associated endocarditis in much of the world \cite{HammondHaley2023}.    

Vancomycin has become the first line antibiotic and the cornerstone of treatment against MRSA since the 1970s, and it is the antibiotic of choice against MRSA endocarditis \cite{Minter2023}. Vancomycin, a glycopeptide, inhibits cell wall synthesis by preventing polymerization of glycan strands and cross linking between the strands via transpeptidase (PBP) \cite{murray_glycopeptides_2019}. The increased reliance on vancomycin since the 1980s has led to the development of vancomycin resistant forms \cite{mcguinness2017vancomycin}. These resistant strains pose a significant clinical challenge, underscoring the ongoing need for new antibiotics and new approaches to manage antibiotic resistance.  

Since the end of the 20th century newer antibiotics, such as daptomycin and linezolid have been developed to target the various forms of staphylococcal aureus that have become resistant to older antibiotics such as vancomycin.  Linezolid, an oxazolidinone antibiotic, has become one of the second line agents used in the treatment of vancomycin resistant MRSA endocarditis \cite{Cong2020}.  AS a protein synthesis inhibitor, linezolid blocks the formation of the bacterial 70S ribosome complex \cite{Shinabarger1997}. Similar to vancomycin, linezolid can be administered intravenously orally.  

Despite the development of newer antibiotics, such as linezolid, bacterial infective endocarditis continues to be a challenge to treat. A better understanding of the factors that drive treatment failure and disease progression is needed. The challenge of vancomycin resistance will require  multifactorial approaches in order to understand how polymicrobial infections are affected by different antibiotics and how their interaction affects resolution and relapse in bacterial endocarditis.  Pharmacokinetic/pharmacodynamic (PK/PD) modeling offers one such approach, enabling researchers to simulate how mixtures of bacterial populations respond to antibiotic exposures over time and under varying host conditions. Monte Carlo simulation is a specific framework used to study virtual patient populations.  It allows investigators to evaluate how vancomycin-sensitive and vancomycin-resistant MRSA populations behave across a range of clinically relevant treatment protocols. By integrating bacterial dynamics, antibiotic pharmacology, and the host innate immune response within a single modeling framework, this approach provides a means of studying the interplay between these forces in a way that human clinical trials cannot.

Mathematical and PK/PD modeling frameworks have proved useful in characterizing the dynamic interactions between MRSA, antibiotic therapy, and the host immune response[cite]. For example, Mikkaichi et al. [year ?] developed an ordinary differential equation-based model incorporating two phenotypically distinct bacterial subpopulations.  One bacterial group was phenotypically normal, while the other group, persister cells, were antibiotic resistant.  Mikkaichi's model was used to simulate persistent MRSA bacteremia that was undergoing treatment with vancomycin.  Using this ensemble modeling approach, Mikkaichi and colleagues identified decreased immune clearance rate of persister cells as the primary determinant of persistent bacteremia, and demonstrated that a persister-killing strategy was the most effective pharmacological approach to achieve bacterial clearance \cite{Mikkaichi2019}. Complementing Mikkaichi's model with a host modeling framework, Vera-Yunca et al. developed a semi-mechanistic PKPD model characterizing the bactericidal dynamics of vancomycin and linezolid against MRSA. A critical finding was that the maximum bacterial killing rate in vivo was 70–75 percent lower than in vitro estimates for both agents\cite{VeraYunca2025}. Together, these studies highlight the gap between in vitro susceptibility data and in vivo treatment response.  The present study uses a dual-compartment mathematical model to examine how the interaction between vancomycin-sensitive and vancomycin-resistant MRSA subpopulations, guideline-directed antibiotic therapy, the host innate immune response, and the rates of bacterial exchange between the bone reservoir and the bloodstream together determine clinical resolution in a cohort of virtual patients. Monte Carlo sampling propagates inter-patient variability in each of these factors through the model, so that the predicted rate of clinical resolution emerges as a distribution of outcomes across the cohort rather than from a single deterministic trajectory.


\section{Clinical Scenario}
A 62‑year‑old man with a history of type 2 diabetes mellitus presented to the emergency department with a history of not feeling well, progressive fatigue and intermittent fever. He had been discharged from a different hospital several weeks prior after treatment for osteomyelitis of the right ankle \cite{MomoduSavaliya2023}. During the previous hospitalization he had underwent surgical debridement and had been discharged to home on oral antibiotics. On admission he was found to have a non healing wound on the right ankle.  The wound was draining clear blood. Cardiac examination was unremarkable.  Transesophageal echocardiogram (TEE) showed posterior leaflet thickening of the mitral valve consistent with bacterial endocarditis. Laboratory studies revealed neutrophil predominant leucocytosis and elevated C‑reactive protein.   

Presumptive diagnosis of Staphylococcal endocarditis was made, two sets of blood cultures were drawn, and he was started on Vancomycin, 1000 mg every 8 hours. The patient continued to spike a fever up to 39 degrees Celsius during his initial hospital course.  On the third day of hospitalization blood culture results showed the presence of Methicillin-resistant Staphylococcus aureus (MRSA) with a vancomycin MIC of $32~\mathrm{mg/L}$. This report was consistent with Vancomycin resistant MRSA. The laboratory report indicated sensitivity to Linezolid.  Vancomycin was discontinued and the patient was started on 600 mg of Linezolid every 12 hours.

Over the subsequent 48 hours the patient's fever curve defervesced. He was continued on Linezolid to complete a 6-week course for MRSA endocarditis as recommended in Infectious Disease Society of America (IDSA) guidelines.
there are two outcomes to this clinical scenario: resolution and relapse.  At the conclusion of treatment up to 30 percent of patients will relapse necessitating surgical intervention to remove the source of infection.  The outcome of clinical relapse will be studied with the proposed dual-compartment pharmacokinetic pharmacodynamic model.




\maketitle
\section{Computational Model}
\subsection{Two Compartment Model}
The experimental system is composed of a blood compartment and reservoir, which represents bone and is the site of the initial infection. Osteomyelitis develops slowly over a period of several weeks before eventually spreading to the blood compartment. The initial condition of the reservoir assumes inoculum of $100$ and $100$ CFU/mL for sensitive and resistant bacterial species respectively.  The blood compartment is assumed to be free of pathological bacterial species.  Thus S(0) = 0 and r(0) = 0 for sensitive and resistant bacterial species.

\subsection{Blood Compartment}
The blood compartment represents a well-perfused environment with abundant nutrient availability and oxygenation, conditions that support relatively rapid proliferation of \textit{Staphylococcus aureus}. There is no attenuation of bacterial killing by the immune system. The carrying capacity of the blood-compartment represents the maximum bacterial load that can be sustained within the vascular compartment. The growth rates of both species and carry capacity are listed in Table one.


\begin{table}[H]
\centering
\large % Increase the text size within this table
\renewcommand{\arraystretch}{1.35} % Adds vertical breathing room between rows
\caption{Blood Compartment Parameter Values}
\vspace{0.5cm}
\begin{tabular}{l l l l}
\hline
\textbf{Parameter} & \textbf{Units} & \textbf{Nominal Value} & \textbf{Reference} \\
\hline
$\rho_s$ & $h^{-1}$ & 0.72 & [11] \\
$\rho_r$ & $h^{-1}$ & $0.649$ & [13]  \\
$B_{\max}$ & CFU/mL & $6000$ & \cite{nieman2022bacterial} \\
\hline
\end{tabular}
\end{table}

\subsection{Reservoir Compartment}
The reservoir compartment models the osseous environment as the primary infection source, assigning distinct growth, death, and immune-clearance parameters to capture its unique physiology. Unlike well-perfused, nutrient-rich blood, infected bone undergoes localized inflammation and microvascular thrombosis, creating hypoxic, avascular zones (sequestra) with poor immune cell access and severely restricted antibiotic diffusion \cite{MomoduSavaliya2023}. Starved of oxygen and nutrients, \textit{S.~aureus} shifts from rapid proliferation to long-term persistence by forming biofilms, invading host cells, or transitioning into metabolically dormant persisters and Small Colony Variants (SCVs) \cite{Beenken2026,Masters2021}. Consequently, the bacteria exhibit slower growth and lower metabolic activity in bone than in blood, driving chronic persistence and reducing antimicrobial susceptibility.

The reservoir acts as a persistent source of bacteria that seed the blood compartment, thereby
contributing to prolonged infection. Exchange between compartments is governed by first-order rate
constants $k_{r \to b}$ (reservoir to blood) and $k_{b \to r}$ (blood to reservoir), with the rate of bacterial movement out of the reservoir and into the blood compartment five times greater than movement from blood to reservoir.




\begin{equation}
R^{b} + S^{b}
\,
\xrightleftharpoons[f_{r \to b}]{f_{b \to r}}
\,
R^{r} + S^{r}
\end{equation}

\begin{table}[H]
    \centering
    \large %Increase the text size within this table
    \renewcommand{\arraystretch}{1.35} % Adds vertical breathing room between rows
    \caption{Reservoir Compartment Parameter Values}
    \vspace{0.5cm}
    \begin{tabular}{l l l l}
    \hline
    \textbf{Parameter} & \textbf{Units} & \textbf{Nominal Value} & \textbf{Reference} \\
    \hline
    $k_{r \to b}$          & $\text{h}^{-1}$     & $5 \times 10^{-5}$ & \cite{Kostakioti2013} \\
    $k_{b \to r}$          & $\text{h}^{-1}$     & $1 \times 10^{-5}$  & 1/5 of efflux \\
    $\rho_{s}^{r}$         & $\text{h}^{-1}$     & $0.19415$              & Estimated from osteomyelitis growth kinetics \\
    $\rho_{r}^{r}$         & $\text{h}^{-1}$     & $0.1765$            &  9.1\% fitness cost\\
    $B_{\max}^{r}$         & CFU/mL              & $1 \times 10^{4}$   & \cite{liu2025development} \\
    $\phi_{LZD}$           & --                  & $0.30$              & \cite{dryden2011linezolid} \cite{spellberg2012systemic} \\
    $\phi_{VAN}$           & --                  & $0.15$              & \cite{MomoduSavaliya2023} \cite{spellberg2012systemic} \\
    \hline
    \end{tabular}
\end{table}

 

\subsection{Immune Response}
The onset of infection activates the body’s innate immune system. White blood cells (WBCs), particularly neutrophils, which are largely quiescent prior to infection, become activated and migrate to sites of bacterial invasion. Once at the site of infection, these cells release bactericidal substances, including reactive oxygen species and proteolytic enzymes, before undergoing cell death. Pro-inflammatory cytokines and chemokines released at the site of infection stimulate the bone marrow to markedly increase WBC production and mobilization, leading to rapid delivery of immune cells through the circulatory system to infected tissues.

Host immune clearance is represented as a constant, first-order killing
term. All four bacterial subpopulations -- sensitive and resistant, in both blood and
reservoir -- are cleared at a rate proportional to their own density and to
a single rate constant, $k_{immune}$ (baseline $0.243~\mathrm{h^{-1}}$).
The reservoir immune response is scaled by $1/10$ of the immune response seen in the blood compartment.  This tenfold reduction in
reservoir efficacy represents the biofilm created by the Staph aureus infection. 


\subsection{Antibiotic Pharmacokinetics}
After the diagnosis of bacterial endocarditis vancomycin is started after blood cultures are obtained. A one-compartment PK model is used for both Vancomycin and Linezolid to simulate serum concentration profiles.  Vancomycin is approximately 50\% protein-bound in plasma, whereas linezolid is more modestly protein-bound ($\sim$30\%, i.e., $\sim$70\% free fraction), a difference the model accounts for via each drug's free-fraction parameter ($f_u$ = 0.50 for vancomycin, 0.70 for linezolid) when converting total plasma concentration to the pharmacologically active free-drug concentration. The model parameters include the following.\\ 

\begin{flushleft}
$\mathrm{Dose}_{i}$: The amount of drug administered per dose. \\[0.5em]
$C_{i}$: The resulting free-drug concentration driving bacterial kill. \\[0.5em]
$\tau_{i}$ : The time between doses. \\[0.5em]
Duration: The total time the drug is administered. \\[0.5em]
$CL_{i}$: Antibiotic clearance. \\[0.5em]
$V_{d}$: The apparent volume (liters) of distribution.
\end{flushleft}






Between doses the concentration decays first-order with elimination rate
$k_{e,i} = CL_i / V_{d,i}$. At each dosing time $t_k = k\,\tau_i$ the
concentration steps up by $\mathrm{Dose}_i / V_{d,i}$ and then decays again,
so that within the $k$-th dosing interval
\begin{equation}
C_i(t) = C_i(t_k^{+})\, e^{-k_{e,i}(t - t_k)},
\qquad t_k \le t < t_{k+1},
\label{eq:pk}
\end{equation}
with the post-dose value set by the accumulation rule
\begin{equation}
C_i(t_k^{+}) = C_i(t_k^{-}) + \frac{\mathrm{Dose}_i}{V_{d,i}},
\qquad C_i(0^{-}) = 0 ,
\label{eq:pkdose}
\end{equation}
for $i \in \{V, L\}$ (vancomycin, linezolid). Each drug is dosed only within
its own course window; outside that window no further doses are added and the
concentration decays toward zero. Multiplying by the unbound fraction
$f_{u,i}$ gives the free-drug concentration that drives bacterial kill.



\begin{table}[H]
\centering
\large %Increase text size within the table
\renewcommand{\arraystretch}{1.35} % Adds vertical breathing room between rows
\caption{Table of Pharmacokinetic Parameter Values}
\vspace{0.5cm} % Adds 0.5cm of space before caption
\begin{tabular}{l l l l}  % 4 columns
\hline
\textbf{Parameter} & \textbf{Units} & \textbf{Nominal Value} & \textbf{Reference} \\
\hline
$E_{\max}^{V}$ & $h^{-1}$ & 1.7 & \cite{VeraYunca2025}\\
$E_{\max}^{L}$ & $h^{-1}$ & 0.8 & \cite{sandberg2010intra} \cite{Scheerans2014}\\
$EC_{50}^{V}$ & mg/L & 0.25  & \cite{Nielsen2007}\cite{VeraYunca2025}\\
$EC_{50}^{L}$ & mg/L &  0.60  & \cite{Nielsen2007}\cite{Scheerans2014}\cite{sandberg2010intra}\\
$V_d$ & Liters & 50 (Vancomycin), 45 (Linezolid) & \cite{Matzke1986} \cite{dryden2011linezolid} \\
Vancomycin Clearance & \text{ml/min} & 83 (5.0 L/h) & \cite{Matzke1986} \\
Linezolid Clearance & \text{ml/min} &  100 (6.0 L/h) & \cite{dryden2011linezolid} \\
\hline
\end{tabular}
\end{table}


Standard values for $V_d$, $K_e$, dosages and dosage time intervals are used during the simulation \cite{vancopk_review}. Normal kidney function is assumed throughout the simulation.  Antimicrobial efficacy in the bone compartment was down-scaled to represent the limited translocation of antibiotics from the blood into osseous tissue.  The exact degree of bone penetration of either antibiotic is not well characterized\cite{koch2022mysteries}.  Bone penetration of Linezolid and Vancomycin were set at 30 and 15 percent respectively.  Therapeutic bone concentrations of antibiotics were attained in the reservoir/bone compartment. A standard protocol for the treatment of infectious endocarditis is used, based on the recommendations of the Infectious Disease Society of America (\cite{Minter2023}). Vancomycin and Linezolid clearances were set at 83 \text{ml/min} (5.0 L/h) and 100 \text{ml/min} (6.0 L/h), respectively.


During simulation the steady state peak Vancomycin concentration reached 36.3~mg/L and a trough concentration of 16.3~mg/L.  The total-drug
area under the concentration-time curve over 24 hours of
$\mathrm{AUC}_{24,\mathrm{ss}} = 600~\mathrm{mg\cdot h/L}$,
independent of volume of distribution or half-life. This falls at the upper
edge of the standard clinical target of 400--600~mg$\cdot$h/L for MRSA
bacteremia and endocarditis\cite{rybak2020therapeutic}. Accounting for the drug's free fraction
($f_u = 0.50$), the corresponding free-drug exposure driving bacterial kill
in the model is $\mathrm{AUC}_{24,\mathrm{ss}}^{\text{free}} = 300~\mathrm{mg\cdot h/L}$.  Linezolid steady-state total-drug trough concentration was calculated at 3.4~mg/L, which is within the recommended range of 2 to 8 mg/L if the MIC is less than 2 mg/L\cite{alsultan2019determining}. 


\subsection{Bacterial System and ODE Model}
The dynamical framework consists of four state variables 
($S$, $R$, $S^r$, $R^r$) governed by a coupled system of four ordinary differential equations. Equations~\eqref{eq:dSdt} and \eqref{eq:dRdt} relate to processes in the vascular compartment, Equations~\eqref{eq:dSrdt} and \eqref{eq:dRrdt} to those in the reservoir compartment.  The natural death rate of bacteria was omitted from the model due to it's small size compared to the other variables.


\begin{equation}
\label{eq:dSdt}
\frac{dS}{dt}
= \rho_{s}\, S \left( 1 - \frac{S + R}{K} \right)
- k_{\text{imm}}\, S
- E^{V}(C_V)\, S
- E^{L}(C_L)\, S
+ f_{r\to b}\, S^r - f_{b\to r}\, S .
\end{equation}


\begin{equation}
\frac{dR}{dt}
= \rho_{r}\,R \left( 1 - \frac{S + R}{K} \right)
- k_{\text{imm}}\, R
- E^{L}(C_L)\, R
+ f_{r\to b}\, R^r - f_{b\to r}\, R .
\label{eq:dRdt}
\end{equation}



\begin{equation}
\frac{dS^r}{dt}
= \rho_{s}^{r} S^{r} \left( 1 - \frac{S^r + R^r}{K^r} \right)
- 0.1\, k_{\text{imm}}\, S^r
- E^{V}(\phi_{V} C_V)\, S^r
- E^{L,r}(\phi_{L} C_L)\, S^r
- f_{r\to b}\, S^r + f_{b\to r}\, S .
\label{eq:dSrdt}
\end{equation}


\begin{equation}
\frac{dR^r}{dt}
= \rho_{r}^{r} R^{r} \left( 1 - \frac{S^r + R^r}{K^r} \right)
- 0.1\, k_{\text{imm}}\, R^r
- E^{L,r}(\phi_{L} C_L)\, R^r
- f_{r\to b}\, R^r + f_{b\to r}\, R .
\label{eq:dRrdt}
\end{equation}



Each drug effect term is a sigmoidal $E_{max}$ relationship with a Hill
coefficient of unity, evaluated at the free-drug concentration reaching
that compartment:
\begin{equation}
E^{i}(C) = \frac{E_{\max}^{i}\, C}{EC_{50}^{i} + C},
\qquad i \in \{V, L\},
\label{eq:hill}
\end{equation}
where $C_V$ and $C_L$ are the free vancomycin and linezolid
concentrations in blood from Eq.~\eqref{eq:pk}, and the reservoir sees
only the penetrated fractions $\phi_V C_V$ and $\phi_L C_L$. The reservoir
linezolid ceiling is additionally scaled to the reservoir growth rate,
$E_{\max}^{L,r} = E_{\max}^{L}\,(\rho_s^r/\rho_s)$, reflecting the reduced
metabolic activity of \textit{S.~aureus} in bone; $E^{L,r}$ denotes the
resulting effect term.

Both compartments have distinct logistic growth dynamics and carrying
capacities, and within each compartment the two strains \emph{share} that
capacity: the logistic term depends on the combined density $S+R$ (blood)
or $S^r+R^r$ (reservoir), so the strains compete directly for the same
niche. Vancomycin acts on the sensitive strain only; linezolid acts on
both. The immune response in the reservoir is scaled by 0.1 of that in the
blood compartment. As detailed previously, the reservoir acts as a seeding
source for both sensitive and resistant populations entering the blood
compartment, and both strains exchange between compartments in both
directions.



\maketitle
\section{Monte Carlo Simulation}

\subsection{Monte Carlo of Linezolid Potency}


To assess how variation in linezolid potency affects predicted
\textit{Staph aureus} bacterial kinetics, $EC_{50,L}$ was sampled from
a log-normal distribution with mean $0.6~\mathrm{mg/L}$ and coefficient
of variation $\mathrm{CV} = 1.1$, with the log-scale parameters
chosen so that the distribution's mean, rather than its median, equals
the baseline value. This value was chosen based another MRSA PKPD model for pneumonia by Vera-Yunca et. al. All remaining parameters were held at baseline.
For each of $n = 1000$ independently sampled $EC_{50,L}$ values, the
full dual-reservoir model was integrated from blood-sterile, reservoir-seeded
initial conditions ($S_b(0){=}R_b(0){=}0$, $S_{res}(0){=}R_{res}(0){=}100$
CFU/mL) over the complete 81-day simulated course (21 days of
pre-treatment equilibration, a 4-day vancomycin course, a consecutive
42-day linezolid course, and a 14-day post-treatment follow-up)
yielding a distribution of trajectories for all four bacterial
subpopulations.

\subsection{Monte Carlo of Linezolid Maximum Effect}

To assess how variation in linezolid's maximum killing capacity affects
predicted \textit{Staph aureus} bacterial kinetics, $E_{max,L}$ was sampled
from a log-normal distribution with mean $1.085~\mathrm{h}^{-1}$ and
coefficient of variation $\mathrm{CV} = 0.307$, with the log-scale
parameters chosen so that the distribution's mean, rather than its median,
equals the baseline value. The mean value for the $E_{max, L}$ was also taken from the prior study from Vera-Yunca and colleagues.  All remaining parameters were held at baseline.
For each of $n = 1000$ independently sampled $E_{max,L}$ values, the full
dual-reservoir model was integrated over the complete 81-day treatment and
follow-up course, yielding a distribution of trajectories for all four
bacterial subpopulations.

Unlike $EC_{50,L}$, lower values of $E_{max,L}$ represent a less potent
drug and shift simulated outcomes toward resistant-strain escape. A
companion deterministic threshold sweep across $E_{max,L} \in
[0.0, 1.0]~\mathrm{h}^{-1}$ was used to locate the precise escape boundary
via bisection, following the same procedure used for $EC_{50,L}$ and
$k_{immune}$.

\subsection{Monte Carlo of host immune response}

To assess how variation in host immune clearance affects predicted
\textit{S. aureus} bacterial kinetics, $k_{immune}$ was sampled from a
log-normal distribution with mean $0.243~\mathrm{h}^{-1}$ and coefficient of
variation $\mathrm{CV} = 0.25$, with the log-scale parameters chosen so
that the distribution's mean, rather than its median, equals the baseline
value. All remaining parameters were held at baseline. For each of
$n = 1000$ independently sampled $k_{immune}$ values, the full
dual-reservoir model was integrated over the complete 81-day treatment and
follow-up course, yielding a distribution of trajectories for all four
bacterial subpopulations.

Lower values of $k_{immune}$ indicate a weaker host immune response,
moving the bacterial system outcomes toward resistant-strain persistence
-- the same direction as $E_{max,L}$, and the opposite of $EC_{50,L}$. A
companion deterministic threshold sweep across $k_{immune} \in
[0.02, 0.22]~\mathrm{h}^{-1}$ was used to locate the precise escape
boundary via bisection, following the same procedure used for $EC_{50,L}$
and $E_{max,L}$.

\subsection{Monte Carlo of Inter-compartmental Transfer Rates}
Unlike the prior three analyses, inter-compartmental transfer rates were examined using a deterministic, log-spaced grid. An initial log-normal Monte Carlo pass ($\mathrm{CV} \approx 112\%$) showed negligible sensitivity of any outcome to either rate across roughly two orders of magnitude around baseline, with curvature emerging only in the extreme upper tail; a grid was therefore used to resolve that region directly.

Each rate was swept independently, one at a time, across a shared grid spanning $1\times10^{-7}$--$1\times10^{0}\,\mathrm{h}^{-1}$, while the other rate was held fixed at its baseline value ($f_{r \to b} = 5\times10^{-5}\,\mathrm{h}^{-1}$; $f_{b \to r} = 1\times10^{-5}\,\mathrm{h}^{-1}$). The grid combined 40 points spaced evenly across the full range with 160 additional points concentrated in the relapse-to-clearance transition band between $1\times10^{-3}$ and $1\times10^{-2}\,\mathrm{h}^{-1}$, and a further 120 across $1\times10^{-2}$--$1\times10^{0}\,\mathrm{h}^{-1}$ to resolve the high-rate regime. A total of 302 points were used in each sweep (\Cref{fig:mc_transfer_rate_sweep_grid}). All other parameters were held fixed. The simulated course spanned 81 days total (1944~h): 21-day pretreatment, 4-day vancomycin, 42-day linezolid, 14-day follow-up. Densities below the $10\,\mathrm{CFU/mL}$ limit of detection (LOD) were treated as cleared.

For each sweep we report three complementary views in Results: the full time-course kinetics as the median and 5th--95th percentile envelope across the grid; four scalar outcomes (peak $S_b$, peak $R_b$, peak $S_{res}$, and final $R_b$) across the full grid; and the same four outcomes at higher resolution within the transition band.

\begin{figure}[H]
    \centering
    \includegraphics[width=1.0\linewidth]{figures/transfer_rate_sweep_grid.png}
    \caption{Log-spaced grid coverage: the 302 rate values actually tested in each sweep, plotted on a shared log axis, with each sweep's baseline value marked for reference. This figure shows sampling density only, not simulation outcomes.}
    \label{fig:mc_transfer_rate_sweep_grid}
\end{figure}





%---------------------------------------------------------


\section{Results}
\subsection{Pre-Treatment Period}
Prior to initiation of Vancomycin, bacterial populations established a quasi-steady-state equilibrium across both compartments. The vancomycin-sensitive \textit{S. aureus} strain in the blood compartment exhibited a growth rate of $\rho_s = 0.72\text{ h}^{-1}$ while the resistant strain growth was $\rho_r = 0.649\text{ h}^{-1}$. In the bone reservoir, local tissue conditions reduced bacterial replication rates to $\rho_s^r = 0.19415\text{ h}^{-1}$ for the sensitive strain and $\rho_r^r = 0.1765\text{ h}^{-1}$ for the resistant strain. Asymmetric exchange between compartments leads to seeding of the blood compartment from the reservoir.

The clinical simulation captures the transition from osteomyelitis to secondary bacteremia and subsequent infective endocarditis.  Vancomycin administration (1000~mg every 8 hours), models clinically recommended steady-state peak and trough plasma concentrations. Vancomycin therapy drives rapid killing of susceptible \textit{S. aureus} subpopulations ($S_b$ and $S_{\text{res}}$).  However, vancomycin exerted no growth inhibition on resistant strains ($R_b$ and $R_{\text{res}}$; MIC = 32~mg/dL).

\subsection{Vancomycin Sensitive Bacteremia}

The 4-day vancomycin course (\texttt{van\_duration} $= 96$ h, starting at \texttt{vanco\_start} $=
504$ h) drives the vancomycin-sensitive strain to extinction in
both blood ($S_b$) and reservoir ($S_{res}$) compartments before
the initiation of linezolid.

Numerically integrating \texttt{dual\_reservoir\_model} with baseline
parameters ($\rho_S = 0.72$, $\rho_R = 0.649$, $\rho_{res,S} = 0.19415$,
$\rho_{res,R} = 0.1765$, $E_{max,V} = 1.7$, $EC_{50,V} = 0.25$,
$B_{max,blood} = 6000$, $B_{max,res} = 10^4$, $f_{r\to b} = 5\times10^{-5}$,
$f_{b\to r} = 1\times10^{-5}$, $Y_0 = [0,0,100,100]$) gives the sensitive-strain
time course in Table~\ref{tab:sb}.

\begin{table}[H]
\centering
\large %Increase text size within the table
\renewcommand{\arraystretch}{1.35} % Adds vertical breathing room between rows
\caption{$S_b$, $R_b$, and $S_{res}$ across the 4-day vancomycin course,
Peak $S_b$ prior
to treatment onset at $5.00\times10^{3}$ CFU/mL (vancomycin
has not yet started).}
\vspace{0.5cm} % Adds 0.5cm of space before caption
\begin{tabular}{l l l l}  % 4 columns
\hline
\textbf{Time} & \textbf{$S_b$ (CFU/mL)} & \textbf{$R_b$ (CFU/mL)} & \textbf{$S_{res}$ (CFU/mL)} \\
\hline
Vancomycin start ($t=504$ h) & $4.43\times10^{3}$ & $1.12$ & $6.16\times10^{3}$ \\
$+24$ h & $<10^{-13}$ & $4.86$ & $<10^{-10}$ \\
$+48$ h & $<10^{-17}$ & $4.74\times10^{3}$ & $<10^{-18}$ \\
$+72$ h & $<10^{-17}$ & $4.89\times10^{3}$ & $<10^{-18}$ \\
Vancomycin end ($+96$ h, $t=600$ h) & $<10^{-17}$ & $4.54\times10^{3}$ & $<10^{-18}$ \\
\hline
\end{tabular}
\label{tab:sb}
\end{table}

By $+24$ h into the course, $S_b$ and $S_{res}$ has fallen below
$10^{-10}$ CFU/mL and
continue to decay toward numerical zero for the rest of the course. Sensitive strain is cleared  within the first day of treatment. Vancomycin's growth inhibition in
blood (Section 2) drives $S_b$ down immediately, with the reservoir clearing in parallel.

\subsection{Vancomycin Resistant Bacteremia}
Following laboratory confirmation of vancomycin-resistant MRSA, therapy was switched to Linezolid (600~mg every 12 hours), yielding recommended steady-state serum concentrations. Linezolid administration drove a reduction in resistant bacterial counts across both compartments. However, due to reduced local immune efficacy and the presence of biofilms MRSA clearance from the reservoir lagged significantly behind vascular clearance.

Separate Monte Carlo simulations were performed on four of the model parameters: $EC_{50,L}$, $E_{max,L}$, $k_{immune}$ and the inter-compartmental transfer rates.  Separate coefficients of variation were used to quantify the uncertainty around baseline estimates. All samples were propagated through the full 81-day dual-reservoir model to quantify how uncertainty in each parameter translates into heterogeneity in clinical outcome.

\subsubsection{$EC_{50}$}

Values for $EC_{50,L}$ are sampled from a log-normal distribution (mean $0.60~\mathrm{mg/L}$ and $\mathrm{CV}=1.1$) for all 1000 virtual patients in the cohort.  \Cref{fig:mc_ec50_lzd_sweep_kinetics.png} shows the full time-course kinetics on the monte carlo simulation
across all trajectories. Both sensitive subpopulations ($S_b$, $S_{res}$) rise to their pretreatment plateau and clear during vancomycin treatment. $R_{res}$ increases during the pre-treatment period and after antibiotic treatment has transitioned to linezolid shows a steady decline.  The $R_{res}$ settles to a reduced-but-detectable plateau for the remainder of therapy. Both compartments return to their pretreatment/carrying-capacity levels within days of linezolid discontinuation in 44\% of virtual patient simulations. The bifurcation of the system from clinical response to relapse occurred at an $EC_{50, L} \approx0.49~\mathrm{mg/L}$ (\Cref{fig:ec50_lzd_threshold_sweep.png}). Increasing values of $EC_{50, L}$ represent decreasing potency of Linezolid. 

\begin{figure}[H]
    \centering
    \includegraphics[width=1.0\linewidth]{figures/ec50_lzd_threshold_sweep.png}
    \caption{Deterministic threshold sweep for $EC_{50,L}$: final $R_b$ (blood)
    and $R_{res}$ (reservoir) at end of follow-up across the swept range, with
    the escape zone shaded and the bisected threshold
    ($EC_{50,L} = 0.49~\mathrm{mg/L}$) marked. Values at or below the limit of
    detection are shown at a display floor so the curves remain continuous.}
    \label{fig:ec50_lzd_threshold_sweep.png}
\end{figure}



\begin{figure}[H]
    \centering
    \includegraphics[width=1.0\linewidth]{figures/mc_ec50_lzd_sweep_kinetics.png}
    \caption{Monte Carlo sweep of $EC_{50, L}$ ($n = 1000$): full
    time-course kinetics for $S_b$, $R_b$, $S_{res}$, and $R_{res}$, shown
    as individual trajectories (light), median (solid), and 5th--95th
    percentile envelope (shaded).}
    \label{fig:mc_ec50_lzd_sweep_kinetics.png}
\end{figure}
The bifurcation point was located by a companion deterministic threshold sweep.

\subsubsection{$E_{\mathrm{max}}$}

Sampling $E_{\mathrm{max},L}$ from a log-normal distribution (mean $1.085~\mathrm{h}^{-1}$ and $\mathrm{CV} = 0.307$) produced a reemergence of bacteremia in 23.8\% (238/1000) of virtual patients. 
\Cref{fig:mc_emax_lzd_sweep_kinetics} illustrates the full time course kinetics across all trajectories from the beginning to the end of the simulation. Values of $E_{\mathrm{max},L} greater than \approx 0.84~\mathrm{h}^{-1}$ lead to clinical resolution of the infection in both compartments. Both the blood compartment ($R_b$) and the reservoir ($R_{\mathrm{res}}$) relapsed (reemergence of bacteremia) below this value, consistent with a weakening of the effectiveness of Linezolid. This bifurcation point was calculated by a deterministic threshold sweep.  

\begin{figure}[H]
    \centering
    \includegraphics[width=1.0\linewidth]{figures/emax_lzd_threshold_sweep.png}
    \caption{Deterministic threshold sweep for $E_{max,L}$: final $R_b$ (blood)
    and $R_{res}$ (reservoir) at end of follow-up across the swept range, with
    the escape zone shaded and the bisected threshold
    ($E_{max,L} = 0.84~\mathrm{h}^{-1}$) marked. Note the direction is reversed
    relative to $EC_{50,L}$: escape occurs \emph{below} the threshold.}
    \label{fig:emax_lzd_threshold_sweep.png}
\end{figure}


\begin{figure}[H]
    \centering
    \includegraphics[width=1.0\linewidth]{figures/mc_emax_lzd_sweep_kinetics.png}
    \caption{Monte Carlo sweep of $E_{\mathrm{max}, L}$: $S_b$, $R_b$, $S_{\mathrm{res}}$ and $R_{\mathrm{res}}$, full time course kinetics shown as individual trajectories (light), median (solid) and 5th-95th percentile envelope (shaded)}
    \label{fig:mc_emax_lzd_sweep_kinetics}
\end{figure}


\subsubsection{Immune Response}

Sampling $k_{immune}$ from a log-normal distribution with mean
$0.243~\mathrm{h}^{-1}$ and $\mathrm{CV} = 0.25$ produced a reemergence of bacteremia in 26.0\% (260/1000) of virtual patients and clinical resolution in the remaining 74.0\% of patients.  The immune response, $k_{immune} \approx 0.2~\mathrm{h}^{-1}$, was the transition point between these two clinical outcomes.  \Cref{fig:mc_immune_response} shows the full time-course kinetics
as the median and 5th--95th percentile envelope across all sampled
trajectories. As described previously, both sensitive subpopulations ($S_b$, $S_{res}$) rise to their
pretreatment plateau and clear during vancomycin treatment. During the course of Linezolid treatment the bacteremia resolves. There is a steady decrease in $R_{res}$ during treatment, but not eradication below the Limit of Detection.  Both compartments return to their
pretreatment/carrying-capacity levels within days of linezolid discontinuation in 26.0\% of virtual patients.  The boundary between these two clinical outcomes was identified by a deterministic threshold sweep at $k_{immune} \approx 0.2~\mathrm{h}^{-1}$.


\begin{figure}[H]
    \centering
    \includegraphics[width=1.0\linewidth]{figures/mc_immune_response.png}
    \caption{Monte Carlo sweep of $k_{immune}$ ($n = 1000$): full
    time-course kinetics for $S_b$, $R_b$, $S_{res}$, and $R_{res}$, shown
    as individual trajectories (light), median (solid), and 5th--95th
    percentile envelope (shaded).}
    \label{fig:mc_immune_response}
\end{figure}

\subsubsection{Sensitivity to compartment exchange rates}
\label{sec:results-transfer-rates}

Two independent deterministic sweeps were performed to evaluate the relationship between the bacteremia and the exchange rates ($f_{b \to r}$ and  $f_{r \to b}$) between the two compartments.  

\begin{figure}[H]
    \centering
    \includegraphics[width=1.0\linewidth]{figures/frb_sweep_kinetics.png}
    \caption{Full time-course kinetics, Sweep 1: $f_{r \to b}$ (reservoir
    $\to$ blood) varied across the log-spaced grid, $f_{b \to r}$ held at
    baseline ($1\times10^{-5}\,\mathrm{h}^{-1}$). Lines show the median and
    5th--95th percentile envelope across all 302 grid points.}
    \label{fig:frb_sweep_kinetics}
\end{figure}

\begin{figure}[H]
    \centering
    \includegraphics[width=1.0\linewidth]{figures/fbr_sweep_kinetics.png}
    \caption{Full time-course kinetics, Sweep 2: $f_{b \to r}$ (blood $\to$
    reservoir) varied across the log-spaced grid, $f_{r \to b}$ held at
    baseline ($5\times10^{-5}\,\mathrm{h}^{-1}$). Lines show the median and
    5th--95th percentile envelope across all 302 grid points.}
    \label{fig:fbr_sweep_kinetics}
\end{figure}

Varying $f_{r \to b}$ revealed two independent transition thresholds bracketing a broad plateau. Below $f_{r \to b} \approx
2.2\times10^{-5}~\mathrm{h}^{-1}$, there is insufficient transfer from reservoir to blood to establish a bacteremia for the time of the simulation.  The peak $S_b$, peak $R_b$, and final $R_b$ remain at zero.  
The initial threshold establishes a plateau of parameter values where peak and final $R_{b}$ are equal.  Linezolid initiation begins a monotonic decline in $R_{b}$ and is extinct at approximately days 55 to 60.  During the time on the plateau the $R_{b}$ remains below the level of detection. As $f_{r \to b}$ rises across the plateau, $R_{res}$ at end-of-course falls steadily and time to rebound shortens as shown in Table~\ref{tab:frb_plateau_relapse}. As the difference between the two transfer rates increases more reservoir bacteria are funneled into the blood compartment.  Once the transfer rate reaches $8.3\times10^{-3}~\mathrm{h}^{-1}$, the transfer becomes uni-directional the final outcome changes from relapse to complete resolution of bacteremia after the linezolid treatment finishes.  The bacteria in the blood are exposed to both immune killing and to Linezolid's full pharmacologic effect. Both $R_b$ and $R_{res}$ MRSA are eliminated together.  Peak $S_{res}$ varies by less than 6\% across the lower six decades of the swept range ($10^{-7}$--$10^{-2}~\mathrm{h}^{-1}$), confirming that reservoir dynamics are essentially insensitive to $f_{r \to b}$ except through this one-directional drainage effect at the extreme
upper end.

\begin{table}[H]
    \centering
    \begin{tabular}{lccccc}
        \hline
        $f_{r \to b}$ (h$^{-1}$) & Peak $R_b$ & Min $R_b$ on LZD
            & $R_{res}$ at LZD end & Final $R_b$ & Rebound lag (d) \\
        \hline
        $3.35\times10^{-5}$          & $4.89\times10^{3}$ & $2.12\times10^{-2}$ & $5.95\times10^{2}$ & $4.89\times10^{3}$ & 3.0 \\
        $5.00\times10^{-5}$\,$^{*}$  & $4.89\times10^{3}$ & $3.17\times10^{-2}$ & $5.93\times10^{2}$ & $4.89\times10^{3}$ & 2.2 \\
        $9.39\times10^{-5}$          & $4.89\times10^{3}$ & $5.92\times10^{-2}$ & $5.91\times10^{2}$ & $4.89\times10^{3}$ & 1.7 \\
        $2.64\times10^{-4}$          & $4.90\times10^{3}$ & $1.63\times10^{-1}$ & $5.80\times10^{2}$ & $4.90\times10^{3}$ & 1.2 \\
        $7.40\times10^{-4}$          & $4.90\times10^{3}$ & $4.36\times10^{-1}$ & $5.51\times10^{2}$ & $4.90\times10^{3}$ & 0.9 \\
        $2.08\times10^{-3}$          & $4.93\times10^{3}$ & $1.05\times10^{0}$  & $4.68\times10^{2}$ & $4.93\times10^{3}$ & 0.7 \\
        $5.83\times10^{-3}$          & $4.99\times10^{3}$ & $1.45\times10^{0}$  & $2.29\times10^{2}$ & $4.99\times10^{3}$ & 0.6 \\
        \hline
    \end{tabular}
    \caption{Within-course relapse dynamics at seven $f_{r \to b}$ values
    spanning the relapse plateau ($f_{b \to r}$ held at baseline
    $1\times10^{-5}~\mathrm{h}^{-1}$). All burdens in CFU/mL. ``Min $R_b$ on
    LZD'' is the minimum blood resistant burden over the 42-day linezolid
    course (days 25--67); ``$R_{res}$ at LZD end'' is the reservoir resistant
    burden at day 67; ``rebound lag'' is the interval between the end of the
    linezolid course and the first time $R_b$ crosses back above the LOD
    (10~CFU/mL). Follow-up ends at day 81, so every relapse listed completes
    inside the 14-day post-treatment window. Trajectories were integrated at
    finer time resolution than the sweep grid (4000 points over 1944~h,
    $\mathrm{rtol}=10^{-9}$, $\mathrm{atol}=10^{-11}$) to resolve the rebound
    crossing. $^{*}$Model baseline rate.}
    \label{tab:frb_plateau_relapse}
\end{table}




Varying $f_{b \to r}$ produced a markedly different, and asymmetric,
picture. Peak $S_b$, peak $R_b$, and final $R_b$ remain flat to within 1\%
across roughly the lower four and a half decades of the sampled range
($10^{-7}$--$10^{-3.5}~\mathrm{h}^{-1}$), consistent with $f_{b \to r}$
acting only as an efflux from blood rather than a seeding route for either
compartment. Peak $S_{res}$, however, is not flat over the upper range: it
rises monotonically from its baseline value ($\sim$6721~CFU/mL) to a maximum
of $\sim$1.17$\times10^{4}$~CFU/mL ($+75\%$) near
$f_{b \to r} = 3\times10^{-2}~\mathrm{h}^{-1}$, as blood-to-reservoir efflux
progressively transfers the sensitive population into the reservoir. Above
that rate the trend reverses and peak $S_{res}$ returns to baseline, because
blood can no longer sustain a population to transfer: peak $S_b$ collapses
from $\sim$4.8$\times10^{3}$~CFU/mL at $3\times10^{-2}~\mathrm{h}^{-1}$ to
below 3~CFU/mL at $5\times10^{-2}~\mathrm{h}^{-1}$, the blood compartment
being drained faster than it can establish. This upper-range behaviour lies
far outside the physiologically plausible interval and leaves the relapse
outcome unchanged across the entire range over which blood infection
establishes at all.

The $f_{b\to r}$ has minimal to no effect on blood compartment establishment or the relapse outcome.  Resistant-strain relapse is robust to
compartment exchange rates across nearly the entire physiologically
plausible range. The reservoir only loses its protective role when
reservoir-to-blood efflux becomes fast enough, and sufficiently
one-directional, to funnel resistant bacteria into blood faster than the
reservoir is replenished.



\section{Discussion}

In this research study we constructed a dual-compartment computational PKPD model to simulate how infective endocarditis evolves over time in a virtual cohort of patients. Across all three pharmacodynamic and host parameters examined, outcome was governed by a sharp threshold rather than a graded dose-response: each parameter separated complete clinical resolution from post-treatment relapse at a single bifurcation point, and the fraction of the cohort relapsing was set by how far that threshold sat from the parameter's calibrated mean. The vascular compartment was seeded by an osteomyelitis secondary to a wound infection in a diabetic patient.  This scenario frequently presents a clinical challenge characterized by frequent
treatment failure. Clinical response in our cohort used standard antibiotic treatment based on recommended guidelines. The study evaluated the effect of $EC_{50}$, $E_{max}$, the innate immune response, and the rates of inter compartmental transfer on the final outcome of the infection. Where available, parameter values were chosen based on prior research and generally accepted population averages.  

Separate monte carlo simulations were used to study how pharmacological parameters, $EC_{50,L}$ and $E_{max,L}$, affect the resolution of bacteremia in the virtual patient cohort.  Prior to the monte carlo analysis the $EC_{50,L}$ transition point separating clinical resolution from relapse was calculated at  $EC_{50,L} \approx 0.49~\mathrm{mg/L}$.  The $EC_{50,L}$ simulation resulted in 56.4\% of the virtual cohort of patients resolving the bacteremia (43.6\% relapse).   The transition point of the $E_{max,L}$ separating resolution from relapse was measured at $0.84~\mathrm{h}^{-1}$.  In the $E_{max,L}$ monte carlo 76.2\% of virtual patients had a resolution of their bacteremia while 23.8\% had a relapse.  


A separate monte carlo simulation was performed on the innate immune response to evaluate how variation in strength affects clinical resolution. The  transition point of the immune response was calculated at $k_{immune} \approx
0.20~\mathrm{h}^{-1}$.  Using the default distribution, 74.0\% of virtual patients had a resolution of bacteremia (26.0\% relapsed).  The simulation illustrates that a significant number of patients, plausibly those with weaker
baseline immune clearance, remain at risk of clinical relapse. This 
quantifies how presence of comorbidities, such as diabetes, immunosuppression and immunosenescence, frequently predict poor patient outcomes.

The extent of clinical resolution was determined by the separation between each parameter's distributional mean and its own calculated threshold, expressed as a $z$-score (Table~S1). A mean lying far from its threshold pushes nearly the entire sampled population to one side (escape fraction near 0\% or 100\%), whereas a mean lying close to it splits the cohort between resolution and relapse. Across all three parameters the $z$-score alone, with no other information about the underlying dynamics, reproduces the simulated escape fraction to within two percentage points. This relationship carries an important qualification for how the reported percentages should be read. The sampling means for $E_{max,L}$ and $k_{immune}$ were not taken unmodified from the literature: each was tuned so that the resulting cohort reproduced the $20$--$30\%$ relapse rate observed clinically after guideline-directed therapy for MRSA endocarditis. The resulting relapse fractions of $23.8\%$ and $26.0\%$ are therefore a consequence of that calibration rather than an independent prediction of it, and the same is true of any model whose free parameters are anchored to an observed outcome rate. What the model does predict independently is the \emph{structure} of the response: that each parameter separates resolution from relapse at a single sharp threshold rather than through a graded dose-response, where those thresholds lie, and how steeply outcome changes across them. That steepness is itself a finding. Because all three calibrated means sit within one standard deviation of their thresholds, modest shifts in any of them move a substantial fraction of the cohort across the boundary: reducing $EC_{50,L}$ by $6.7\%$, for example, halves the predicted $k_{immune}$ relapse fraction. The clinically meaningful quantity is therefore the narrowness of the band separating the two outcomes, not the precise percentage falling either side of it.


The exchange kinetics between the two compartments affected the
relapse phenotype asymmetrically.  The $f_{r \to b}$ (reservoir$\to$blood) exhibits two distinct transitions (\Cref{fig:frb_sweep_kinetics}). The lower transition, below which  bacteremia is not established and a second transition above which the resistant bacteremia is fully cleared after treatment is complete. The $f_{b \to r}$ (blood$\to$reservoir), in contrast, does not impact resolution or relapse of the bacteremia.


Whether a coupling parameter's timescale matters depends on its topology:
whether it is a net seeding route into a compartment that can re-establish
infection, or a reversible efflux into a compartment the cells can freely
return from. $f_{r \to b}$ gates reservoir-to-blood reseeding directly, so
its timescale relative to the linezolid course determines whether the
reservoir can reseed blood before treatment ends. $f_{b \to r}$ only
redistributes bacteria into a compartment they already occupy and can
reflux back from, so however its timescale lines up with therapy duration,
it cannot by itself flip the relapse-versus-clearance outcome. A coupling
parameter is therefore worth measuring precisely only when (1) its
reciprocal falls within the range of the therapy's own phase durations,
\emph{and} (2) it lies on a net seeding path rather than a reversible
efflux. Given how difficult exchange and penetration rates are to measure
\emph{in vivo} in bone and cardiac vegetations, this distinction tells the
modeler not only which parameters need precise measurement, but which
topological role makes that precision worth pursuing.


These computational results help to mathematically define the well
understood clinical challenges in the management of Staphylococcal
infective endocarditis. Follow-up blood cultures used to assess
microbiological clearance may provide false reassurance of complete
resolution, as vascular sterilization can occur while substantial bacterial
burdens persist in avascular or biofilmed bone tissue. Clinical outcomes
are dependent on the individual patient's innate immune system,
compartmental coupling, and the pharmacodynamic strength of the antibiotics
used in treatment.

\section{Limitations and Future Directions}

This simulation study carries several limitations driven by current gaps
in empirical data.
 First, the model relies on several assumptions about the bone reservoir
compartment. Key parameters --- inter-compartmental transfer rates,
intra-bone bacterial growth and death rates, and the baseline densities of
sensitive and resistant populations --- are not directly measured and had
to be estimated. The transfer rate between bone and blood was modeled as
asymmetric, with efflux from bone into blood set five-fold greater than
influx in the opposite direction. While this directional asymmetry is
biologically plausible (a growing, seeded focus draining into the
bloodstream faster than blood re-seeds it), the specific five-fold
magnitude is an estimated value, not one derived from measured data.

Second, the host immune response is represented by a single first-order
clearance rate, $k_{immune}$, applied uniformly to each compartment (with
a lower effective rate in bone than in blood, reflecting reduced immune
access to a deep-seated focus). This rate had to be estimated, and the
model does not attempt to represent the underlying cellular or molecular
mechanisms of clearance: neutrophil recruitment and killing dynamics, and
the networks of cytokines\footnote{Signaling molecules that regulate
inflammation, cell growth, and immune cell activity.} and
chemokines\footnote{A family of small proteins that attract and guide
immune cells to specific sites of infection.} that coordinate them, are
omitted entirely in favor of this single aggregate rate.

Third, the blood carrying capacity for \textit{S. aureus} is based on a
limited set of measurements using PCR-based quantification of bacterial
DNA, and the maximal carrying capacity used in the model was derived from
this same limited dataset.

Fourth, published estimates of $EC_{50}$ do not consistently distinguish
between total drug concentration and the free (unbound) fraction in
plasma; this model used free serum concentration throughout its
pharmacokinetic component, which may not match the basis of every $EC_{50}$
value drawn from the literature.

Fifth, the resistant strain carries a fitness cost in both compartments
--- approximately $9.9\%$ in blood ($\rho_r = 0.649$ versus
$\rho_s = 0.72~\mathrm{h}^{-1}$) and $9.1\%$ in the reservoir
($\rho_r^r = 0.1765$ versus $\rho_s^r = 0.19415~\mathrm{h}^{-1}$) --- but
the reservoir value is fixed by calibration rather than by measurement,
and is tightly constrained on both sides. It sits just above the rate
below which the resistant reservoir population fails to persist at all,
and only slightly below the rate at which the resistant strain wins the
pre-treatment competition for blood outright, leaving no visible
sensitive-strain bacteremia. The reservoir growth rates should therefore
be read as calibrated quantities chosen to reproduce the clinical
scenario, not as independently measured biological values.

A further limitation concerns the calibration itself. The sampling means for
$E_{max,L}$ and $k_{immune}$ were adjusted so that the simulated cohort
reproduced the clinically reported $20$--$30\%$ relapse rate, which means the
reported relapse fractions are anchored to that target rather than derived
independently of it. Two consequences follow. First, the agreement between the
model's relapse rate and the clinical one is not evidence that the model is
correct, since it was imposed. Second, because the calibrated means end up
within one standard deviation of their thresholds, the escape fractions are
sensitive to small parameter changes --- a $6.7\%$ reduction in $EC_{50,L}$
halves the $k_{immune}$ relapse fraction, and a $0.01$ change in the linezolid
reservoir penetration fraction moves it across most of its range. Reported
percentages should therefore be treated as indicative of regime rather than as
precise risk estimates, and the threshold locations --- which are re-derived by
root-finding and do not depend on the sampling distributions --- are the more
robust output.

Finally, the in vitro bacterial doubling times used for calibration may
not fully reflect the metabolically dormant or small-colony-variant (SCV)
states typical of chronic biofilm infections, nor do they
necessarily reflect the actual replication rate of MRSA within the
circulatory system in vivo. All growth-rate calibration in this model is
ultimately anchored to in vitro measurements, and in vivo replication
dynamics -- particularly within a biofilm or vegetative focus -- remain a
source of uncertainty this model cannot resolve.

Additional limitations stem from constraints on available patient-level
data and the scope of the sensitivity analyses performed. Because
patient-level data on model parameters such as growth rates of
MRSA in blood pharmacodynamic and other physiological variables remain limited, the
model could not be formally validated against clinical outcome data. The
Monte Carlo sensitivity analyses were also univariate: each sweep varied
a single parameter ($EC_{50,L}$, $E_{max,L}$, or $k_{immune}$) while
holding every other parameter fixed at baseline. Real
inter-patient variability is likely joint and correlated.  Impaired renal function, can simultaneously reduce drug clearance and also alter immune status.  This univariate approach cannot capture
these interaction effects between parameters that vary together in an actual
patient population. Similarly, no attempt was made to introduce
inter-patient pharmacokinetic variability: volume of distribution,
clearance, and protein binding were held fixed at population-average
values for every simulated patient. Finally, the simulation
used a single fixed treatment protocol throughout, with no exploration
of missed doses, alternative dosing regimens, or source control (e.g. 
surgical debridement or hardware removal), all of which materially
affect outcomes in real bone and endocarditis infections.



\bibliographystyle{ieeetr}
\bibliography{references}






\end{document}
